% Chapter 1, Introduction.
% Generated from Dissertation/Chapter1_Introduction.md by tools/md2tex.py.
%% The draft-status line, the provenance notes and the [edit] markers in the
%% Markdown source are not part of the submitted text and are absent here.
%% Markdown is the source. Regenerate this file when it changes; do not hand-edit.

\chapter{Introduction}
\label{ch:intro}
\lhead{Chapter 1. \emph{Introduction}}

\section{Background and motivation}
\label{sec:background}

Between 2017 and 2022, \citet{OShaughnessy2023Commercial} reports, the United Kingdom fell from fourth to tenth in the world for the number of commercial clinical trials initiated, and enrolments in commercially sponsored NIHR studies fell by 44 per cent, from roughly 50,000 patients a year to 28,000. The decline matters because patients lose access to treatments available only through trial participation.

The Department of Health and Social Care introduced a national 150-day metric, from regulatory application to first participant recruitment, to improve the United Kingdom's competitiveness in commercial clinical research \citep{anon2025Support}. At site level, each trust has 90 days from regulatory approval or site selection, whichever is later: 60 days to site opening and a further 30 days to recruit its first participant. This dissertation examines that 60-day set-up window.

At Lancashire Teaching Hospitals (LTHTR), Research and Innovation (R\&I) records the milestone dates, reports performance against the national metric, and decides which trials to take on and when each is passed to the departments that must prepare them. Pharmacy is one of those departments, and no site opens until it has issued its green light.

Trial set-up follows different pharmacy pathways, including routine review, aseptic preparation and electronic prescribing, so the clinical trials pharmacy team's capacity for about two new commercial trials a month is not interchangeable across studies. Demand is projected to grow by 10 to 20 per cent, increasing the importance of making the existing pathway capacity visible and usable.

Departments prepare their parts of the set-up concurrently, so pharmacy holding a 60-day goal of its own is not a problem in itself. The difficulty is that site opening waits for the last required track to finish, so a pharmacy overrun transfers directly to the site-opening date. On the available Capability and Confirmation Agreement proxy, only 3 of 24 trials fell within the 60-day allowance between 2024 and 2026; the median interval was 144 days. The tracker holds no site-confirmation date, so this later agreement stands in for it (Section~\ref{sec:inpractice}).

\section{Operational problem}
\label{sec:opproblem}

Meeting the site-opening target at LTHTR cannot be addressed by pharmacy alone. Set-up capacity is fixed and distributed across general and specialist teams, while R\&I decides which studies pharmacy receives and when, and therefore the volume, mix and timing of work entering the pharmacy pathway. Progressing one study consumes capacity that may be needed by another, yet the feasible number and specialist mix of studies are not visible at portfolio level. The structural problem is therefore to select the studies that LTHTR can progress within existing pharmacy capacity and identify the capacity constraint limiting further set-up. Chapter 3 describes this arrangement in full, and sets out which part of it this study models.

\section{Research objective}
\label{sec:objective}

The objective of this study is to determine how much of LTHTR's existing pharmacy capacity can be converted into completed set-ups of commercial medicine trials, and where effort is stranded when it cannot. Staffing is held constant throughout.

The dissertation does not model clinical site set-up, governance and committee calendars, or staffing expansion, which lies outside the commissioned scope; Section 3.2 defines the operational and evidential boundaries of these exclusions.

The dissertation develops a repeatable planning model that identifies which trials can be committed within the 60-day window, distinguishing a trial's individual requirements from the outcome achievable by the portfolio as a whole. It selects the plan with the most on-time set-ups and, among tied plans, the least weighted lateness across the queue.

The contribution is the translation itself: an optimisation model of LTHTR's own pharmacy set-up process, its capacity, its pathways and its shared approval step, calibrated on the trust's records, which makes the question of how many trials can be set up on time answerable rather than a matter of opinion.

\section{Research questions}
\label{sec:questions}

Stated in terms the study can measure, the objective becomes a single critical question.

\textbf{RQ. How many commercial medicine trials can LTHTR commit to pharmacy set-up within its 60-day set-up window using the capacity it already has, and what limits that number?}

It is therefore approached through three sub-questions, each answered by one output of the planning model (see Table~\ref{tab:questions}).

\begin{table}[htbp]
  \centering
  \caption{The critical research question and its three sub-questions*}
  \label{tab:questions}
  \begin{tabularx}{\textwidth}{@{}LL@{}}
    \toprule
    \textbf{Focus} & \textbf{Sub-question} \\
    \midrule
    Capacity ceiling & SQ1: How many commercial medicine trials, and which of them, can LTHTR commit to pharmacy set-up within the number of set-ups it can take on in a cycle? \\
    \addlinespace
    Binding constraint & SQ2: Which constraint binds on-time completion across the observed decision points, and what changes when the cycle commitment limit is varied? \\
    \addlinespace
    Selection rule and its robustness & SQ3: Does selecting on how many trials finish inside 60 days, with waiting time breaking ties, outperform the arrival order used today, and does any advantage survive when arrival timing, trial mix and set-up duration are varied in simulation? \\
    \bottomrule
  \end{tabularx}
\end{table}

\section{Dissertation structure}
\label{sec:structure}

The chapters follow the argument from question to answer. Chapter 2 identifies the gap in the measurement, selection and queueing literature, and Chapter 3 states the pharmacy set-up problem and formalises it as a planning model. Chapter 4 sets out the data, the planning inputs, the experimental design and the assumptions each rests on, and Chapter 5 reports what the model returns on the trust's own queues and under simulated variation. Chapter 6 answers the research questions, compares the findings with published work, and states the limitations and what they imply for practice and further research.

---

